\section*{Block sensitivity and randomized communication (author source pp.12--13)}
A block $B\subseteq[n]$ is sensitive for a function $f$ at $x$ iff $f(x)\ne f(x^B)$. The block sensitivity $\operatorname{bs}(f,x)$ of $f$ at $x$ is the maximum number of disjoint sensitive blocks for $f$ at $x$. The block sensitivity $\operatorname{bs}(f)$ of $f$ is the maximum sensitivity $\operatorname{bs}(f,x)$ over all points $x\in\bits^n$.
We use the notation $x^B$ for $x\in\bits^n$ and $B\subseteq[n]$ to denote $x$ with all the coordinates in $B$ flipped.
In a randomized communication game, Alice and Bob have access to an unlimited amount of random bits, and they can use it to decide which bit to send next. We say that a (private coin) randomized communication protocol $\Pi$ computes a function $f:X\times Y\to Z$ with error $\varepsilon$ if
\[\Pr_{r_A,r_B}[\Pi(x,y,r_A,r_B)=f(x,y)]\ge1-\varepsilon,\qquad\forall x\in X,y\in Y,\]
where $r_A$ and $r_B$ are the strings of random bits used by Alice and Bob, respectively. The cost of a randomized protocol is the maximum number of bits that can be sent. We denote by $R_\varepsilon(f)$ the minimal cost of a private coin randomized communication protocol that computes $f$ with error $\varepsilon$. The randomized communication complexity of $f$ is defined as $R(f):=R_{1/3}(f)$.
\section*{6.1 Reductions in communication complexity (author source p.35)}
Consider a pair of two-party functions $f_i:X_i\times Y_i\to\bits$ for $i=1,2$. We say that $f_1$ reduces to $f_2$, denoted $f_1\le f_2$, if the communication matrix of $f_1$ is a submatrix of the communication matrix of $f_2$. Equivalently, $f_1\le f_2$ iff there exist one-to-one mappings $\pi_A$ and $\pi_B$ such that
\[f_1(x,y)=f_2(\pi_A(x),\pi_B(y)),\qquad\forall(x,y)\in X_1\times Y_1.\]
